\section{Evaluation Protocol}
\label{sec:evaluation}

\subsection{Three Groups and Fixed Correct Answers}

For each concept, we evaluated target questions, neighboring-topic
questions, and questions from unrelated topics. The target group
measures change on the concept to be erased; the other groups assess
preservation at different degrees of topical proximity. Each group
contained 50 selection and 50 test questions. The target and
neighboring questions followed EMBER's existing splits
\citep{suslik2026ember}; we sampled the unrelated questions from
other topics.

We used Claude in a dedicated session to rewrite the multiple-choice
questions as declarative sentence stems while retaining their correct
answers. For example, ``Which river is associated with the founding
of Rome?'' became ``The river associated with the founding of Rome
is'', with ``The Tiber River'' as the correct continuation. This
format avoided relying on the 1B base models to produce a gradeable
response to a multiple-choice instruction, which they often failed
to do in preliminary tests.

Ranking the supplied options would give disproportionate weight to
three particular wrong answers, although many other incorrect
continuations are possible. We therefore scored only the correct
continuation.

\subsection{Answer NLL Difference}

For question $i$, let $x_i$ be the sentence stem and
$y_{i,1},\ldots,y_{i,n_i}$ the tokens of its correct continuation.
We measured model $M$'s mean negative log-likelihood (NLL) over
those answer tokens:
\begin{equation}
\ell_M(i)
= -\frac{1}{n_i}\sum_{t=1}^{n_i}
\log p_M(y_{i,t}\mid x_i,y_{i,<t}).
\end{equation}
Each token was scored given the stem and the preceding correct-answer
tokens; the model did not generate a response.

To compare $M$ with reference model $R$, either the full model or the
concept-excluded twin, we computed the paired difference
$\Delta_{M,R}(i)=\ell_M(i)-\ell_R(i)$. A positive difference means
that $M$ assigned less probability to the correct answer than $R$.

\subsection{Teacher-Forced Full-Vocabulary KL}

Answer NLL measures support for the supplied answer. To measure
broader changes in prediction, we also compared the models' full
next-token distributions at each answer position. Let $p_{M,t}$
denote model $M$'s distribution over the vocabulary given the stem
and the preceding correct-answer tokens. For concept-excluded twin
$T$ and comparison model $M$, we computed
\begin{equation}
\operatorname{KL}_{T\to M}(i)
= \frac{1}{n_i}\sum_{t=1}^{n_i}
D_{\mathrm{KL}}(p_{T,t}\,\|\,p_{M,t}).
\end{equation}

We kept the twin on the left when comparing it with both the full
model and each erased model. This choice matters because KL is
asymmetric: reversing its arguments can change the values and even
the ordering of models. Using the twin as the common reference asks
how closely each model reproduces the twin's predictions, with
deviations weighted by the twin's probability distribution. Lower
KL indicates greater similarity at the evaluated positions.

The full vocabulary includes incorrect next tokens, so KL can detect
changes among alternatives even when the correct token's probability
changes little. Because the models are evaluated along the correct
continuation, KL does not measure their predictions after an
incorrect token is chosen.

We report the mean and standard deviation across the 50 held-out
test questions in each target, neighboring, and unrelated group.
For NLL, these summarize paired question-level differences; for KL,
they summarize question-level divergence.